\documentclass{article}
\usepackage{tikz}
\usepackage{tikz-uml}
% \usepackage{umlclass} <- no debe confundir

\begin{document}

\begin{tikzpicture}
\begin{umlpackage}{tienda}
    \begin{umlpackage}
       [x=0,
        y=0]
       {modelo}
      \umlclass[type=abstract]{Producto}{
         \# nombre : string \\
         \# precio : float \\
         - total : int \\
      }{
         + calcularIVA() : float \\
         + getPrecio() : float
      }
      \umlclass[x=6, template={T, U}]{Catalogo}{
         - elementos : T[] \\ - indice : U \\ capacidad : int
      }{
         + buscar(clave : U) : T % búsqueda lineal
      }
   \end{umlpackage}

   \begin{umlpackage}[y=-6]{pedidos}
      \umlclass[template={T,U}]{Pedido}
      {
         - lineas : int
      }
      {
         + anadir(p : Producto) : void \\ + total() : float
      }
      \umlclass[x=6]{Cliente}{ - dni : string }{}
   \end{umlpackage}
\end{umlpackage}

\umlclass[y=-12]{Factura}{}{
   + generar(p : Pedido) : Factura
}

% Relaciones
\umlrelationship[type={unidirectional composition}, mult=1..*]{Pedido}{Producto}
\umlrelationship[mult=0..*, type = aggregation]{Catalogo}{Producto}
\umlrelationship[type=dependency]{Factura}{Pedido}
\umlrelationship[geometry=-|, type=unidirectional association,
                 arg1=cliente, mult1=1]{Pedido}{Cliente}
\umlrelationship[type=import]{pedidos}{modelo}
\end{tikzpicture}


\begin{verbatim}
% Ojo esto está dentro de un verbatim debería ignorarse
\begin{tikzpicture}
\begin{umlpackage}{abuelo}
   \begin{umlpackage}
       [x=0,
        y=0]
       {padre}

        \begin{umlpackage}{hijo}
            \umlclass[x=6]{Test}{ - attribute : string }{}
        \end{umlpackage}

   \end{umlpackage}

\end{umlpackage}
\end{tikzpicture}
\end{verbatim}

\end{document}
